\lyseis{2}{Πρόσθεση, αφαίρεση, πολλαπλασιασμός: οι ιδιότητες που κρύβονται}
\sol{1} Στο $13+5=18$: προσθετέοι και άθροισμα. Στο $20-8=12$: μειωτέος, αφαιρετέος και διαφορά. Στο $7\cdot 6=42$: παράγοντες και γινόμενο.
\sol{2} α) αντιμεταθετική της πρόσθεσης \quad β) προσεταιριστική της πρόσθεσης \quad γ) επιμεριστική του πολλαπλασιασμού ως προς την πρόσθεση \quad δ) το 1 δεν αλλάζει το γινόμενο \quad ε) ο πολλαπλασιασμός με το 0 δίνει 0 \quad στ) το 0 δεν αλλάζει το άθροισμα.
\sol{3} α) 5.200 \quad β) 370 \quad γ) 4.900.000 \quad δ) 8.000.
\sol{4} α) $40-17=23$ \quad β) $90-26=64$ \quad γ) $32+15=47$ \quad δ) $80-35=45$. Έλεγχος, π.χ. στο δ: $80-45=35$.
\sol{5} α) $4\cdot 25\cdot 17=100\cdot 17=1.700$ \quad β) $8\cdot 125\cdot 39=1.000\cdot 39=39.000$ \quad γ) $(12+88)+(47+53)=100+100=200$ \quad δ) $5\cdot 2\cdot 73=10\cdot 73=730$.
\sol{6} α) $89\cdot(7+3)=89\cdot 10=890$ \quad β) $(23+77)\cdot 49=100\cdot 49=4.900$ \quad γ) $76\cdot(13-3)=76\cdot 10=760$ \quad δ) $284\cdot(100-1)=28.400-284=28.116$ \quad ε) $35\cdot(100+2)=3.500+70=3.570$.
\sol{7} α) Λάθος: το $4-9$ δεν υπάρχει στους φυσικούς. β) Λάθος: $(20-8)-2=10$, ενώ $20-(8-2)=14$. γ) Λάθος: $5\cdot(2+3)=25$, ενώ $5\cdot 2+3=13$. δ) Λάθος: $7\cdot 0=0$. ε) Σωστό. στ) Σωστό.
\sol{8} α) 190 \quad β) $(122+78)+25=225$ \quad γ) 462 \quad δ) 2.726 (έλεγχος: $4.595+2.726=7.321$) \quad ε) $60-16=44$.
\sol{9} $1+2+3+4=10$, \ $1+2+34=37$, \ $12+34=46$, \ $1\cdot 2\cdot 3\cdot 4=24$.
\sol{10} $98\cdot 5=(100-2)\cdot 5=500-10=490$ τοστ.
\sol{11} Αφετηρία ως γέφυρα: $47$ χλμ. Γέφυρα ως βενζινάδικο: $112-47=65$ χλμ. Βενζινάδικο ως χωριό: $175-112=63$ χλμ. Χωριό ως θάλασσα: $230-175=55$ χλμ. Βενζινάδικο ως θάλασσα: $230-112=118$ χλμ. Έλεγχος: $63+55=118$ και $47+65+63+55=230$.
\sol{12} Μειωτέος είναι τα αυτοκόλλητα που είχε: $Μ=17+38=55$. Έλεγχος: $55-17=38$.
\sol{13} Στην αφαίρεση $82-47$ το 2 στις μονάδες είναι μικρότερο από το 7. Δεν επιτρέπεται να γυρίσουμε τους αριθμούς ανάποδα: δανειζόμαστε μία δεκάδα και κάνουμε $12-7=5$, και μετά $7-4=3$. Σωστό: 35. Έλεγχος με $Μ=Α+\Delta$: $47+35=82$. Με το 45 του Άρη βγαίνει $47+45=92$, που δεν είναι 82.
\sol{14} $48\cdot 25=12\cdot 4\cdot 25=12\cdot 100=1.200$ πλακάκια.
\sol{15} $13\cdot 8=(10+3)\cdot 8=10\cdot 8+3\cdot 8=80+24=104$ τ.μ. Ένα ορθογώνιο $10\times 8$ και ένα $3\times 8$, κολλημένα.
\sol{16} $3\cdot 99=3\cdot(100-1)=300-3=297$ €.
\sol{17} Τα ζευγάρια είναι 50 και το καθένα κάνει 101, άρα $50\cdot 101=5.050$. Για το $1+\dots+10$: 5 ζευγάρια $(1+10),(2+9),\dots$ και καθένα κάνει 11, άρα $5\cdot 11=55$. Έλεγχος με πρόσθεση: $1+2+3+4+5+6+7+8+9+10=55$.
\sol{18} Οι αριθμοί είναι 7 και 5: $7+5=12$ και $7-5=2$.
\sol{19} Τέσσερις σακούλες με 0 καραμέλες έχουν συνολικά $0+0+0+0=0$ καραμέλες. Γενικά $4\cdot 0=0+0+0+0=0$.
\sol{20} Το ορθογώνιο με 3 σειρές των 5 έχει 15 τετράγωνα. Αν το γυρίσεις στο πλάι, βλέπεις 5 σειρές των 3, και τα τετράγωνα είναι πάλι 15. Η εικόνα είναι διαφορετική (3 ομάδες των 5 ή 5 ομάδες των 3), το αποτέλεσμα ίδιο.
\sol{21} $(10-5)-2=5-2=3$ και $10-(5-2)=10-3=7$. Οι τιμές διαφέρουν, άρα η αφαίρεση δεν είναι προσεταιριστική: έχει σημασία ποιες δύο αφαιρέσεις γίνονται πρώτα.
\sol{22} $3\cdot 3=9$, $9\cdot 3=27$, $27\cdot 3=81$, $81\cdot 3=243$, $243\cdot 3=729$. Έχει 3 ψηφία. Πιο σύντομα: $3^6$, «το τρία στην έκτη».
